\documentclass[11pt]{article}
\usepackage[T1]{fontenc}
\usepackage[utf8]{inputenc}
\usepackage{lmodern,microtype}
\usepackage[a4paper,margin=28mm]{geometry}
\usepackage{amsmath,amssymb,amsthm,mathtools}
\usepackage{xcolor}
\usepackage[colorlinks=true,linkcolor=blue!45!black,citecolor=blue!45!black,urlcolor=blue!45!black]{hyperref}
\usepackage{enumitem}
\setlist{leftmargin=*,itemsep=2pt}
\setlength{\emergencystretch}{2em}
\newtheorem{theorem}{Theorem}[section]
\newtheorem{lemma}[theorem]{Lemma}
\newtheorem{proposition}[theorem]{Proposition}
\newtheorem{corollary}[theorem]{Corollary}
\theoremstyle{definition}\newtheorem{definition}[theorem]{Definition}
\theoremstyle{remark}\newtheorem{remark}[theorem]{Remark}
\DeclareMathOperator{\Hom}{Hom}\DeclareMathOperator{\Ext}{Ext}\DeclareMathOperator{\Tor}{Tor}
\DeclareMathOperator{\Der}{Der}\DeclareMathOperator{\Spec}{Spec}\DeclareMathOperator{\gr}{gr}
\newcommand{\workbenchurl}{https://github.com/KokunoYumeto/schemes-and-moduli-workbench}
\newcommand{\recordurl}{https://github.com/KokunoYumeto/mathematics-workbenches}
\title{Cohen lifting for Grothendieck's adic algebras\\and Deligne's algebraic extensions}
\author{}
\date{23 November 2026}
\begin{document}
\maketitle
\begin{abstract}
Grothendieck's study of formally smooth topological algebras includes a lifting assertion withdrawn from EGA because its compactness argument was incorrect, while Deligne's correspondence asks for presentations of algebraic extensions that remain flat after arbitrary Artinian lifts of the equations. We prove the withdrawn lifting assertion by applying linear compactness only to square-zero kernels. We also prove existence and uniqueness up to isomorphism of flat, complete, formally smooth lifts of noetherian adic algebras across square-zero extensions of the base, without a locality hypothesis, and construct the presentations asked for by Deligne in every cardinality. Three proposed corrections concerning meromorphic functions and generic formal fibres are already in Kleiman's work or EGA's errata and are cited as such. The mathematics and text were produced by GPT-6 Astra (ChatGPT Pro mode), with the machine checks described at the end; no human has checked them.
\end{abstract}

\section{Grothendieck's lifting question and the statements}

The point of the first result is that the withdrawn assertion of EGA remains true. In EGA $0_{\mathrm{IV}}$, Proposition 19.3.12, pp.~85--86, Grothendieck considers a formally smooth topological $A$-algebra $B$ and an inverse limit $C$ of Artinian rings. For a closed topologically nilpotent ideal $J$ of $C$, he states that
\begin{quote}
\guillemotleft\ tout $A$-homomorphisme continu $u:B\to C/\mathfrak J$ se factorise en $B\xrightarrow{v}C\to C/\mathfrak J$\guillemotright.
\end{quote}
The erratum in EGA IV, fourth part, Err$_{\mathrm{IV}}$~8, p.~347, deletes the proposition because the spaces of lifts used in its proof are not affine linear spaces \cite{EGAIV1,EGAIV4}. This objection concerns the proof. The first theorem below restores the statement with an explicit proof and a slightly more flexible nilpotence hypothesis.

All rings in this paper are commutative with identity. A topological ring has a basis of neighborhoods of zero consisting of ideals. A topological $A$-algebra $B$ is \emph{formally smooth} if every continuous $A$-algebra map from $B$ to a discrete quotient $D/K$, where $K^2=0$, lifts to a continuous map to $D$. Iteration gives the same property for nilpotent kernels. An \emph{adic algebra} in the results below is separated and complete for the powers of a finitely generated ideal.

\begin{theorem}[Grothendieck's lifting assertion]\label{thm:proartinian}
Let $B$ be a formally smooth topological $A$-algebra. Let $C$ be a topological $A$-algebra with a directed basis of open ideals $U_\lambda$ such that every $C/U_\lambda$ is Artinian and
\[
 C\simeq\varprojlim_\lambda C/U_\lambda
\]
as topological rings. Suppose $J\subset C$ is closed and its image in every $C/U_\lambda$ is nilpotent. If the square of every open ideal of $B$ is open, then every continuous $A$-algebra map $B\to C/J$ lifts continuously to $C$.
\end{theorem}

The inverse system need not have a countable cofinal subsystem. This is where ordinary successive lifting through a sequence of quotients is insufficient. We first solve the compatibility problem with a square-zero kernel, where the equations for the difference of two lifts are linear. Closed powers of $J$ then give a countable tower to which that result applies.

The second result concerns the complete noetherian adic algebras of Grothendieck's Cohen-algebra dossier. Definition~4.7 is on cote~22, sheet~12. On sheet~13, after reduction to a square-zero ideal, he writes:
\begin{quote}
\guillemotleft\ On obtient alors un probl\`eme de nature essentiellement infinit\'esimale, dont nous ignorons s'il existe toujours une solution\guillemotright.
\end{quote}
Here the sheet numbers are the archive's pencil numbering, also printed in the digital footer \cite{Archive}. Unlike the complete-local lifting theorem in EGA $0_{\mathrm{IV}}$, 19.7.2, pp.~108--109, the statement below makes no locality assumption. We state the topological module conditions explicitly. An $R$-module $M$ with a linear topology is \emph{formally projective} if every continuous $R$-linear map from $M$ to a discrete quotient module lifts through the quotient. It is \emph{strictly formally projective} if it has a basis of open submodules $V$ for which $M/V$ is projective. These are the conditions of EGA $0_{\mathrm{IV}}$, 19.2.1 and 19.2.3 \cite{EGAIV1}.

\begin{theorem}[Noetherian adic Cohen lifting]\label{thm:adic}
Let $A\twoheadrightarrow A_0$ be a surjection of noetherian rings with square-zero kernel $I$. Let $B_0$ be a noetherian $A_0$-algebra, separated and complete for the $J$-adic topology, and formally smooth over the discrete ring $A_0$. There is a noetherian flat $A$-algebra $B$ with a specified isomorphism $B/IB\simeq B_0$. For the inverse image $Q$ of $J$, the ring $B$ is separated, complete, and formally smooth over $A$.

The lift is unique up to an $A$-isomorphism inducing the identity on $B_0$. Its group of such automorphisms is
\[
 \Der_{A_0}(B_0,I\otimes_{A_0}B_0).
\]
If $B_0$ is formally projective, or strictly formally projective, as a topological $A_0$-module, then $B$ has the corresponding property over $A$.
\end{theorem}

Existence of a lift is a deformation problem, whereas formal smoothness is initially a lifting property for maps. The connection is the cotangent complex. Its reductions modulo $J^n$ are projective modules in degree zero. This kills the obstruction group and the group classifying different lifts. The regularity theorem for completed local maps, Popescu's theorem, and Neeman's detection by residue fields are the substantive inputs \cite[15.50.2, 16.12.1]{Stacks}\cite[Lemma 2.12, p.~527]{Neeman}. No finite generation of the projective module of differentials is assumed.

P\'erez Rodr\'iguez proves existence of affine smooth formal lifts in \cite[Definition~1.3 and Corollary~4.3; author version, pp.~3 and~13]{Perez2008}. Her completed-cotangent characterization \cite[Proposition~5.1; author version, p.~21]{Perez2016} also assumes that the morphism is of pseudo finite type, meaning that its reduction modulo ideals of definition is of finite type. Here $B_0/J$ need not be of finite type over $A_0$, so those results do not by themselves give the full statement above.

For a recent finite-type comparison, Kerz and Tamme prove a simplicially enriched lifting equivalence for topologically finite-type formally smooth algebras over a complete principal-adic base \cite[\S3, Proposition~3.1, pp.~3--4]{KerzTamme}. Cohen rings also enter Feld's recent AI-assisted comparison of rational Gersten kernels: a complete mixed-characteristic discrete valuation ring is compared with a Cohen subring and then with Witt vectors \cite[Theorem~1.1, p.~2, and Corollary~4.3, p.~12]{FeldCohen}. That work applies the classical complete-local structure theorem to $K$-theory. Our deformation assertion concerns nonlocal adic algebras over a square-zero extension, and uses neither comparison as an input.

Deligne's algebraic-extension question, in cote~23, sheet~62, asks for equations whose arbitrary Artinian lifts stay flat:
\begin{quote}
\guillemotleft\ si on rel\`eve les $P_i$ \`a un anneau artinien de corps r\'esiduel $k$, soit $A$, $A[T_\alpha]/(\widetilde P_i)$ soit plat sur $A$\guillemotright.
\end{quote}
The following statement includes inseparable extensions and arbitrary cardinalities.

\begin{theorem}[Deligne's algebraic-extension question]\label{thm:triangular}
For every algebraic extension $K/k$ of fields, there is an ordinal $\lambda$ and a presentation
\[
 K=k[T_\alpha\mid\alpha<\lambda]/(P_\alpha\mid\alpha<\lambda)
\]
such that $P_\alpha$ is monic in $T_\alpha$ and uses only variables of index at most $\alpha$. For every Artinian local ring $(A,\mathfrak m)$ with $A/\mathfrak m=k$, and every choice of polynomial lifts
\[
 \widetilde P_\alpha\in A[T_\beta\mid\beta<\lambda],
\]
the quotient by all the $\widetilde P_\alpha$ is a free $A$-module. The lifted equations may involve later variables; they need not remain triangular.
\end{theorem}

The finite triangular construction appears in \cite[Tag~07BV, Lemma~10.158.11]{Stacks}. The theorem makes its arbitrary-cardinality form and its behavior under arbitrary equation lifts explicit. Triangularity is needed only in the residue field. Every finite increasing list of residue equations is regular. Nilpotence of the maximal ideal transfers that regularity to a flat quotient upstairs, and filtered colimits finish the argument.

\subsection*{Earlier corrections and the scope of this paper}

The meromorphic-function errors in EGA IV 20.1.3 were corrected by Kleiman in 1979. His example on p.~203 is the square-zero thickening of a projective line by a negative line bundle; his example on p.~204 shows that a meromorphic stalk can be smaller than the total quotient ring of its local ring \cite{Kleiman}. The definitions and the noetherian repair are also in Stacks, Tags 01X1 and 0EMG. Similarly, EGA III, second part, Errata et addenda, Liste~2, part~A, entry (III, 4.3.6), already replaces $K'=K\otimes_A\widehat A$ in III 4.3.6 by the inclusion $K'\supset K\otimes_A\widehat A$ \cite{EGAIII2}. These are known corrections and are not additional theorems of this paper.

The first theorem makes a precise distinction between a false compactness argument and a true lifting statement. The second gives the complete deformation statement needed for nonlocal adic algebras. The third makes the cardinality and the freedom in lifting the equations explicit. All proofs below specify their hypotheses; none asserts a first occurrence of these conclusions.

\section{Compatible lifts over Artinian quotients}

This section proves Theorem~\ref{thm:proartinian}; the essential point is to use compactness only for affine linear equations.

\begin{lemma}[Linear compactness]\label{lem:compact}
Let $(M_i)_{i\in I}$ be Artinian modules over a ring $R$. Consider affine $R$-linear equations in variables $z_i\in M_i$, each involving finitely many variables. If every finite subsystem has a solution, the whole system has a solution.
\end{lemma}
\begin{proof}
A family of cosets of submodules of one Artinian module has nonempty intersection whenever its finite intersections are nonempty. Choose, by the descending chain condition, a minimal submodule among the submodules underlying the finite intersections. Its coset is contained in every member of the family: further intersection is nonempty and cannot decrease the underlying submodule.

Well-order the variables. Suppose that values of an initial segment have been chosen so that finite solvability persists after adding equations fixing those values. For the next variable, project the solution sets of all finite subsystems to its Artinian module. These are cosets of submodules with the finite intersection property. Choose a point in their intersection. This preserves finite solvability. At a limit stage a finite subsystem uses only finitely many previous choices, so the property persists there as well. The resulting assignment satisfies every equation.
\end{proof}

\begin{lemma}[Closed quotients]\label{lem:quotient}
Under the hypotheses on $C$ in Theorem~\ref{thm:proartinian}, every closed ideal $D$ satisfies
\[
 C/D\simeq\varprojlim_\lambda C/(D+U_\lambda)
\]
as topological rings.
\end{lemma}
\begin{proof}
Closedness says $D=\bigcap_\lambda(D+U_\lambda)$ and gives injectivity. Given compatible elements on the right, choose variables in the Artinian $C$-modules $C/U_\lambda$. Require their images to equal the prescribed elements and require compatibility. These are affine linear equations. A finite collection is solved at a single dominating index by choosing one representative in $C$ and projecting. Lemma~\ref{lem:compact} gives a compatible family and hence an element of $C$. Both topologies have the kernels of the displayed quotient maps as a neighborhood basis.
\end{proof}

\begin{proof}[Proof of Theorem~\ref{thm:proartinian}]
First suppose $J^2=0$. Put $C_\lambda=C/U_\lambda$ and $J_\lambda=(J+U_\lambda)/U_\lambda$. For each $\lambda$ choose one algebra lift $w^0_\lambda:B\to C_\lambda$ of the given map to $C_\lambda/J_\lambda$; formal smoothness permits this. All other lifts are $w^0_\lambda+z_\lambda$, with $z_\lambda:B\to J_\lambda$ satisfying
\[
\begin{aligned}
 z_\lambda(b+b')&=z_\lambda(b)+z_\lambda(b'),\\
 z_\lambda(bb')&=w^0_\lambda(b)z_\lambda(b')+w^0_\lambda(b')z_\lambda(b),\\
 z_\lambda(a)&=0\quad(a\in A),\qquad z_\lambda(1)=0.
\end{aligned}
\]
We regard each $z_\lambda(b)$ as a variable in the Artinian $C$-module $J_\lambda$. The displayed equations and the equations making the lifts compatible are affine linear. Every finite subsystem has a solution: choose an index above those appearing, use formal smoothness there, and project. Lemma~\ref{lem:compact} therefore gives compatible lifts. Their inverse limit is an algebra map $v:B\to C$. By Lemma~\ref{lem:quotient}, its reduction modulo $J$ is the given map.

This map is continuous. Fix $\lambda$. Some open ideal $V$ of $B$ maps to zero in $C_\lambda/J_\lambda$, so its lifted image is in $J_\lambda$ and $V^2$ maps to zero in $C_\lambda$. Since $V^2$ is open, each component of $v$ is continuous.

For arbitrary $J$, let $J_n=\overline{J^n}$. Multiplication and the ideal topology give $J_n^2\subset J_{n+1}$. Lemma~\ref{lem:quotient} makes each $C/J_n$ complete and an inverse limit of Artinian rings. The map $C/J_{n+1}\to C/J_n$ has closed square-zero kernel. Successive applications of the square-zero case produce compatible continuous maps $v_n:B\to C/J_n$.

We check that their limit has target $C$. For every $\lambda$, nilpotence gives $J^{n(\lambda)}\subset U_\lambda$, hence $J_{n(\lambda)}\subset U_\lambda$. Thus $\bigcap_nJ_n=0$. Given compatible classes in $C/J_n$, choose representatives $c_n$. Modulo each $U_\lambda$ they are eventually constant. Their eventual values define an element $c$ of $C$; closedness of $J_n$ gives $c-c_n\in J_n$. It follows that
\[
 C\xrightarrow{\sim}\varprojlim_n C/J_n.
\]
This is a homeomorphism: every $U_\lambda$ contains some $J_n$ and $U_\lambda/J_n$ is open in that quotient. The inverse limit of the maps $v_n$ is the desired continuous lift.
\end{proof}

\begin{remark}
The restriction to linear equations is essential. For an infinite field $k$, the equations $(x-a)y_a=1$, indexed by $a\in k$, are finitely soluble but have no simultaneous solution. In the lifting problem, the product $z(b)z(b')$ disappears exactly when the kernel is square-zero. The passage to closed powers ensures that it disappears at every stage where compactness is invoked.
\end{remark}

\section{The cotangent complex of an adic algebra}

This section supplies the obstruction vanishing needed for Theorem~\ref{thm:adic}. We use cohomological degrees, with $L_{B/R}$ in degrees at most zero. We use the standard base-change and transitivity properties of the cotangent complex \cite[92.6--92.7]{Stacks}.

\begin{lemma}\label{lem:adiccotangent}
Let $R\to B$ be a homomorphism of noetherian rings, where $R$ is discrete and $B$ has the $J$-adic topology. If $B$ is formally smooth over $R$, then, for every $n\geq1$,
\[
 L_{B/R}\otimes_B^{\mathbf L}B/J^n
 \simeq (\Omega_{B/R}\otimes_B B/J^n)[0],
\]
and the module on the right is projective over $B/J^n$.
\end{lemma}
\begin{proof}
Write $B_n=B/J^n$ and $P_n=\Omega_{B/R}\otimes_BB_n$. To prove projectivity, let $M\twoheadrightarrow N$ be a surjection of $B_n$-modules and let $d:B\to N$ be an $R$-derivation. Leibniz's rule gives $d(J^{n+1})=0$. Thus
\[
 B\longrightarrow B_n\ltimes N,\qquad b\longmapsto(\bar b,d(b))
\]
is continuous. Formal smoothness lifts it through $B_n\ltimes M$. Its second coordinate is a derivation lifting $d$. The representing property of $P_n$ proves it is projective.

We next compute the derived fibres. Fix $\mathfrak q\supset J$, put $\mathfrak p=\mathfrak q\cap R$, and write $\ell=\kappa(\mathfrak q)$. The local homomorphism $R_\mathfrak p\to B_\mathfrak q$ is formally smooth for the maximal-ideal topology of its target. A test map from $B_\mathfrak q$ restricts to a $J$-adically continuous map on $B$; a lift on $B$ sends the denominators outside $\mathfrak q$ to units and extends to the localization. If the original map kills a power of $\mathfrak q$, its lift kills twice that power, so it is continuous.

Pass to the maximal-adic completions. Formal smoothness persists under completion \cite[Tag 07ED]{Stacks}. A formally smooth map between complete noetherian local rings is regular \cite[Proposition 15.50.2]{Stacks}. Popescu's theorem expresses a regular noetherian ring map as a filtered colimit of smooth maps \cite[Tag 07GC]{Stacks}. Cotangent complexes commute with filtered colimits, and for a smooth map the cotangent complex is the projective module of differentials in degree zero \cite[Tags 08S9 and 08R5]{Stacks}. Consequently the cotangent complex of a regular noetherian map is a flat module in degree zero.

For a noetherian local ring $(T,\mathfrak m)$, flatness of completion and cotangent base change give
\[
 L_{\widehat T/T}\otimes_{\widehat T}^{\mathbf L}T/\mathfrak m=0.
\]
Using this for $R_\mathfrak p$ and $B_\mathfrak q$, together with transitivity, identifies
\[
 L_{B/R}\otimes_B^{\mathbf L}\ell
 \simeq L_{\widehat B_\mathfrak q/\widehat R_\mathfrak p}
           \otimes_{\widehat B_\mathfrak q}^{\mathbf L}\ell.
\]
It is concentrated in degree zero.

Set $E_n=L_{B/R}\otimes_B^{\mathbf L}B_n$. Its zeroth cohomology is $P_n$. Because $P_n$ is projective, the identity of $P_n$ lifts to a morphism $P_n[0]\to E_n$. Its cone has zero derived tensor product with every residue field of $B_n$: the preceding paragraph gives concentration in degree zero and the morphism is the identity on that degree. Neeman's residue-field detection theorem for the unbounded derived category of a noetherian ring \cite[Lemma 2.12, p.~527]{Neeman} says that this cone vanishes. This proves the assertion.
\end{proof}

\begin{lemma}\label{lem:extvanish}
If $B$ in Lemma~\ref{lem:adiccotangent} is $J$-adically complete and $N$ is a finite $B$-module, then $\Ext_B^i(L_{B/R},N)=0$ for every $i>0$.
\end{lemma}
\begin{proof}
Noetherian completeness gives $N=\varprojlim N/J^nN$. The transition maps are surjective, so this is also the derived inverse limit. Derived Hom in its second variable preserves derived inverse limits. Adjunction and Lemma~\ref{lem:adiccotangent} therefore give
\[
 \mathbf R\Hom_B(L_{B/R},N)
 \simeq \mathbf R\varprojlim_n
 \Hom_{B/J^n}(P_n,N/J^nN)[0].
\]
The transition maps of these Hom groups are surjective. Indeed, a map at stage $n$ is a derivation into $N/J^nN$, and projectivity of $P_{n+1}$ lifts it through $N/J^{n+1}N\twoheadrightarrow N/J^nN$. The inverse system is countable and Mittag--Leffler. Its derived limit has no positive cohomology, which is the required vanishing.
\end{proof}

\section{Existence, uniqueness, and the module conditions}

We now construct the lift and verify each property of Theorem~\ref{thm:adic}.

\begin{proof}[Proof of the algebra assertions of Theorem~\ref{thm:adic}]
First, $B_0$ is flat over $A_0$. Completeness gives $J\subset\operatorname{rad}(B_0)$ by the geometric-series inverse of $1-jb$. At every maximal ideal of $B_0$, the localization is formally smooth for its maximal-ideal topology, by the localization argument in Lemma~\ref{lem:adiccotangent}. The local flatness theorem for formal smoothness \cite[Tag 07NQ]{Stacks} applies, and flatness is local.

Put $N=I\otimes_{A_0}B_0$, a finite $B_0$-module. The obstruction to a square-zero extension
\[
 0\longrightarrow N\longrightarrow B\longrightarrow B_0\longrightarrow0
\]
compatible with $A\twoheadrightarrow A_0$ and the map $i\mapsto i\otimes1$ lies in $\Ext^2_{B_0}(L_{B_0/A_0},N)$ \cite[Lemma 92.16.1]{Stacks}. Lemma~\ref{lem:extvanish} kills this obstruction. For the resulting extension, the ideal $IB$ is exactly $N$. The multiplication map
\[
 I\otimes_A B=I\otimes_{A_0}B_0\longrightarrow B
\]
is therefore injective. Thus $\Tor_1^A(A_0,B)=0$. Together with $A_0$-flatness of $B_0$, this implies $A$-flatness of $B$: first apply change of rings to modules killed by $I$, and then filter any $A$-module $M$ by $IM$ and $M/IM$ to obtain $\Tor_1^A(M,B)=0$.

The ring $B$ is noetherian. For any ideal, its image in $B_0$ is finitely generated, and its intersection with $N$ is finite over $B_0$. Lifts of the first generators together with generators of the intersection generate the ideal in $B$.

Let $Q$ be the inverse image of $J$. On $N$, the $Q$-adic module topology is the $J$-adic topology. It is complete, since $N$ is finite over the complete noetherian ring $B_0$. Exactness of completion on finite modules over a noetherian ring, applied to the displayed extension, gives
\[
 0\longrightarrow N\longrightarrow\widehat B\longrightarrow B_0\longrightarrow0.
\]
The comparison map from the original sequence has identity end maps. Therefore $B\simeq\widehat B$, proving separation and completeness.

To prove formal smoothness, let $M$ be a $B$-module killed by a power of $Q$. Flat base change gives
\[
 L_{B/A}\otimes_B^{\mathbf L}B_0\simeq L_{B_0/A_0}.
\]
The modules $IM$ and $M/IM$ are killed by powers of $J$ as $B_0$-modules. Lemma~\ref{lem:adiccotangent} and adjunction show that $\Ext_B^1(L_{B/A},-)$ vanishes on both. The long exact sequence then gives its vanishing on $M$.

Now take a continuous $A$-map $B\to D/K$ with $D$ discrete and $K^2=0$. Pull back the extension $D\to D/K$ to $B$. Its kernel is a $B$-module killed by a power of $Q$, so the vanishing just proved splits the extension as an $A$-algebra. This gives an algebra lift to $D$. If the original map kills $Q^s$, the lift kills $Q^{2s}$ and is continuous.

The obstruction theorem also gives the classification of lifts and their automorphisms \cite[Lemma 92.16.1]{Stacks}. The respective groups are
\[
 \Ext^1_{B_0}(L_{B_0/A_0},N)\quad\text{and}\quad \Ext^0_{B_0}(L_{B_0/A_0},N).
\] The former vanishes and the latter is $\Der_{A_0}(B_0,N)$. Any isomorphism inducing the identity modulo $IB$ preserves $Q$, hence is continuous. This proves all the algebra assertions.
\end{proof}

\begin{lemma}\label{lem:formalproj}
Let $A$ be noetherian, $I^2=0$, and $B$ a noetherian flat $A$-algebra endowed with the $Q$-adic topology for an ideal $Q$. If $B/IB$ is formally projective over $A/I$ for its quotient topology, then $B$ is formally projective over $A$.
\end{lemma}
\begin{proof}
For a linearly topologized module $V$ with a decreasing basis $V_n$ of open submodules, write
\[
 T_V^i(M)=\varinjlim_n\Ext_R^i(V/V_n,M).
\]
Here $T_V^0(M)$ is the group of continuous maps to the discrete module $M$. The lifting property is equivalent to $T_V^1(M)=0$ for every $M$. One implication follows from the long exact sequence of Ext. For the other, embed $M$ into an injective module and apply the lifting property to the quotient of that injective module.

If $IM=0$, the low-degree change-of-rings sequence gives
\[
\begin{split}
0\to\Ext_{A/I}^1(B/(IB+Q^n),M)
 &\to\Ext_A^1(B/Q^n,M)\\
 &\to\Hom_{A/I}(\Tor_1^A(B/Q^n,A/I),M).
\end{split}
\]
Flatness of $B$ identifies the Tor module with $(IB\cap Q^n)/IQ^n$. Artin--Rees for the finite $B$-submodule $IB\subset B$ gives $c$ with
\[
 IB\cap Q^{n+c}\subset IQ^n.
\]
Thus the transition map in the Tor system is zero after increasing the index by $c$. Its dual direct limit is zero. The direct limit of the left Ext groups is zero by formal projectivity downstairs. Hence $T_B^1(M)=0$ when $IM=0$.

For arbitrary $M$, both $IM$ and $M/IM$ are killed by $I$. The long exact Ext sequence, followed by the exact filtered colimit, now proves $T_B^1(M)=0$. This is formal projectivity of $B$.
\end{proof}

\begin{lemma}\label{lem:strictproj}
In Lemma~\ref{lem:formalproj}, suppose $B$ is complete and separated for the $Q$-adic topology. If $B/IB$ is strictly formally projective over $A/I$, then $B$ is strictly formally projective over $A$.
\end{lemma}
\begin{proof}
Write $B_0=B/IB$ and $J=QB_0$. Choose a decreasing basis $V_n$ of open submodules of $B_0$, beginning with $V_0=B_0$, such that $V_n\subset J^n$ and each $B_0/V_n$ is projective over $A/I$. Split the surjections $B_0/V_{n+1}\to B_0/V_n$, and let $D_{0,n}$ be their kernels. Each $D_{0,n}$ is projective, and completeness gives a topological isomorphism
\[
 B_0\simeq\prod_{n\geq0}D_{0,n},
\]
whose tails map onto $V_n$.

Every projective $A/I$-module lifts to a projective $A$-module, in arbitrary rank. Express it as the image of an idempotent matrix on a free module, lift the column-finite matrix to an endomorphism $u$, and put $\delta=u^2-u$. Since $\delta^2=0$ and $u\delta=\delta u$, the endomorphism
\[
 e=u+(1-2u)\delta
\]
satisfies $e^2=e$. Its image is the required lift. Apply this to obtain projective modules $D_n$ lifting $D_{0,n}$.

The specified map $D_{0,n}\to J^nB_0$ lifts through $Q^n\twoheadrightarrow J^nB_0$, by projectivity of $D_n$. Denote a lift by $f_n:D_n\to Q^n$. Completeness defines a continuous map
\[
 h:P=\prod_{n\geq0}D_n\longrightarrow B,
 \qquad (x_n)\longmapsto\sum_n f_n(x_n).
\]
Since $I$ is finitely generated, $IP=\prod_n ID_n$. The reduction of $h$ is the chosen isomorphism onto $B_0$. Its cokernel is equal to its product with $I$, and is zero because $I^2=0$. If $T$ is the kernel, flatness of $B$ makes $T/IT\to P/IP$ injective. The reduction of $h$ is an isomorphism, so $T=IT=0$. Hence $h$ is an algebraic isomorphism.

We verify that its inverse is continuous. Fix a tail $P_{\geq n}$. Choose $u$ with $J^u$ contained in the reduction of $h(P_{\geq n})$. Artin--Rees gives $r$ such that $IB\cap Q^r\subset IQ^u$. Choose $m\geq\max(n,r)$ and then $s\geq r$ such that $J^s$ lies in the reduction of $h(P_{\geq m})$. For $b\in Q^s$, choose $x\in P_{\geq m}$ with $b-h(x)\in IB$. This difference lies in $Q^r$, hence in $IQ^u$. If $z\in Q^u$, it agrees modulo $IB$ with some $h(y)$, where $y\in P_{\geq n}$. Multiplication by $I$ kills the difference, giving $Iz\subset h(IP_{\geq n})$. Therefore $b\in h(P_{\geq n})$. We have proved $Q^s\subset h(P_{\geq n})$.

The images of the tails are consequently open and form a basis. Their quotients are finite products of projective $A$-modules, hence projective. This proves strict formal projectivity and completes Theorem~\ref{thm:adic}.
\end{proof}

\section{Presentations of algebraic extensions}

This section proves Theorem~\ref{thm:triangular}, keeping track of why arbitrary lifts of the equations are permitted.

\begin{lemma}\label{lem:koszul}
Let $(A,\mathfrak m,k)$ be Artinian local, let $P=A[T_\alpha\mid\alpha\in\Lambda]$, and let $f_1,\ldots,f_s\in P$ have images forming a regular sequence in $\overline P=k[T_\alpha]$. Then $P/(f_1,\ldots,f_s)$ is free over $A$ and the Koszul complex on the $f_i$ has no positive homology.
\end{lemma}
\begin{proof}
Filter the Koszul complex $D_\bullet$ by $\mathfrak m^jD_\bullet$. Its graded piece in filtration degree $j$ is
\[
 (\mathfrak m^j/\mathfrak m^{j+1})\otimes_k
 K_\bullet(\overline P;\overline f_1,\ldots,\overline f_s).
\]
The residue Koszul complex is exact in positive degrees. The filtration is finite, since $\mathfrak m^h=0$ for some $h$. A positive-degree cycle can therefore be corrected by a boundary to raise its filtration, and after at most $h$ corrections it is zero. The same filtration argument in degree zero gives
\[
 \gr_\mathfrak m(P/(f))
 \simeq\gr_\mathfrak m(A)\otimes_k\overline P/(\overline f).
\]
Lift a $k$-basis of the residue quotient to $P/(f)$. The resulting map from a free $A$-module induces the displayed isomorphism on associated graded modules. Induction along the finite filtration proves both its surjectivity and its injectivity. All combinations are finite, even when the chosen basis is infinite.
\end{proof}

\begin{proof}[Proof of Theorem~\ref{thm:triangular}]
Well-order a set of field generators $(x_\alpha)_{\alpha<\lambda}$ of $K/k$. The subalgebra $K_{<\alpha}=k[x_\beta\mid\beta<\alpha]$ is a field: each of its elements belongs to a finite algebraic extension of $k$ contained in that subalgebra. Let $Q_\alpha$ be the monic minimal polynomial of $x_\alpha$ over $K_{<\alpha}$. Represent its finitely many coefficients by polynomials in earlier generators. This produces a monic polynomial $P_\alpha$ using only variables of index at most $\alpha$.

Transfinite induction shows that the quotient by the relations of index less than $\alpha$ is $K_{<\alpha}$. At successor stages this is the minimal-polynomial presentation; at limit stages it is the union of the preceding fields. Hence the full quotient is $K$.

Any finite subset of the relations, ordered by increasing index, is a regular sequence in the full polynomial ring. Before imposing $P_\alpha$, none of the earlier relations contains $T_\alpha$. The quotient is still a polynomial ring in that variable over a possibly large ring. A monic polynomial is a nonzerodivisor in a polynomial ring over any coefficient ring, by comparison of leading coefficients. The successive quotients are nonzero because they map to $K$.

Apply Lemma~\ref{lem:koszul} to an arbitrary finite set of lifted relations. Each such quotient is free, hence flat, over $A$. Their filtered colimit is the quotient by all the lifted relations. Filtered colimits preserve flatness, so it is flat over $A$.

Finally every flat module $M$ over an Artinian local ring is free. Lift a basis of $M/\mathfrak mM$ to obtain a map $F\to M$ from a free module. Its cokernel satisfies $C=\mathfrak mC$, hence is zero by nilpotence. If $L$ is the kernel, flatness makes $L/\mathfrak mL\to F/\mathfrak mF$ injective; the latter maps isomorphically to $M/\mathfrak mM$. Thus $L=\mathfrak mL=0$. This finishes the proof.
\end{proof}

For the smallest inseparable example, take $a\notin k^p$, so $k(a^{1/p})=k[T]/(T^p-a)$. Over $A=k[\epsilon]/(\epsilon^2)$, consider the nonmonic lift $f=T^p-a+\epsilon T^{p+1}$. Direct calculation gives
\[
 f=(1+\epsilon T)(T^p-a+\epsilon aT).
\]
The first factor is a unit with inverse $1-\epsilon T$. The quotient by $f$ is therefore the quotient by the monic second factor, and is free on $1,T,\ldots,T^{p-1}$. The theorem extends this behavior to arbitrary Artinian lifts and to infinite presentations.

\section{Scope of the conclusions}

The three theorems concern arbitrary pro-Artinian targets, complete noetherian adic source algebras, and arbitrary algebraic field extensions, respectively. The noetherian hypotheses in the second theorem enter the regularity comparison, detection of cotangent complexes by residue fields, exactness of completion, and Artin--Rees. This argument does not prove the same assertion for nonnoetherian adic algebras. The third theorem does not require the lifted equations to be triangular and does not assert a regular-sequence presentation for an arbitrary ring extension.

A comparison with the general literature on deformations of nonlocal formal schemes remains necessary to determine whether Theorem~\ref{thm:adic}, with both module conditions, has an earlier identical formulation. No conclusion of the paper depends on treating the known EGA corrections as new.

\section*{Method and records}

The mathematics and text were produced by GPT-6 Astra (ChatGPT Pro mode). Earlier proof checks were recorded by Claude Opus 5.5 at effort max. Separate GPT-6 Astra instances independently re-derived all three main theorems and the elementary meromorphic examples from their statements; these are checks by the same model. The complete proofs and source comparisons here still require human review. No human has checked this text. The text is dedicated to the public domain under CC0 1.0. The workbench is \url{\workbenchurl}; the permanent address of the records is \url{\recordurl}.

\begin{thebibliography}{9}
\bibitem{EGAIV1} A.~Grothendieck, with J.~Dieudonn\'e, \emph{\'El\'ements de g\'eom\'etrie alg\'ebrique IV, premi\`ere partie}, Publ. Math. IHES \textbf{20} (1964), 5--259. Chapter $0_{\mathrm{IV}}$, 19.2, 19.3.12, 19.7.2, 19.8.1.
\bibitem{EGAIV4} A.~Grothendieck, with J.~Dieudonn\'e, \emph{\'El\'ements de g\'eom\'etrie alg\'ebrique IV, quatri\`eme partie}, Publ. Math. IHES \textbf{32} (1967), 5--361. 20.1.3, pp.~227--228; Err$_{\mathrm{IV}}$~8, p.~347.
\bibitem{EGAIII2} A.~Grothendieck, with J.~Dieudonn\'e, \emph{\'El\'ements de g\'eom\'etrie alg\'ebrique III, seconde partie}, Publ. Math. IHES \textbf{17} (1963), 5--91. Errata concerning III 4.3.6.
\bibitem{Kleiman} S.~L.~Kleiman, \emph{Misconceptions about $K_X$}, Enseign. Math. (2) \textbf{25} (1979), 203--206. \url{https://doi.org/10.5169/seals-50379}.
\bibitem{Perez2008} M.~P\'erez Rodr\'iguez, \emph{Basic deformation theory of smooth formal schemes}, J. Pure Appl. Algebra \textbf{212} (2008), 2381--2388. Definition~1.3 and Corollary~4.3. Author version: \url{https://martapr.webs.uvigo.es/Investigacion/deformation-v3.2.pdf}.
\bibitem{Perez2016} M.~P\'erez Rodr\'iguez, \emph{Deformation of formal schemes through local homology}, J. Algebra \textbf{445} (2016), 78--102. Proposition~5.1. Author version: \url{https://arxiv.org/abs/2402.02879}.
\bibitem{KerzTamme} M.~Kerz and G.~Tamme, \emph{A remark on crystalline cohomology}, Math. Z. \textbf{309} (2025), article~48. Proposition~3.1, article p.~4. \url{https://doi.org/10.1007/s00209-024-03675-6}.
\bibitem{FeldCohen} N.~Feld, \emph{Universal test rings for rational Gersten injectivity}, AI-assisted preprint, October 2026, version~2. Theorem~1.1, p.~2; Corollary~4.3, p.~12. \url{https://hexagonmath.org/2610.00619v2}.
\bibitem{Neeman} A.~Neeman, with an appendix by M.~B\"okstedt, \emph{The chromatic tower for $D(R)$}, Topology \textbf{31} (1992), 519--532. Lemma 2.12, p.~527.
\bibitem{Stacks} The Stacks Project Authors, \emph{The Stacks Project}, \url{https://stacks.math.columbia.edu}. Tags 01X1, 0EMG, 07ED, 07NQ, 07PK, 07GC, 08R5, 08S9, 08SM.
\bibitem{Archive} A.~Grothendieck and P.~Deligne, \emph{Archives Grothendieck}, Universit\'e de Montpellier, cotes 22 and 23. \url{https://grothendieck.umontpellier.fr}.
\end{thebibliography}
\end{document}
